% --- ARCHIVO: Capitulos/posicionamiento.tex ---

\chapter{Comienzo de Partida}

\begin{loreblock}
``El terreno nunca es neutral. Antes de que los guerreros crucen espadas, ya hay un vencedor marcado por quien leyó mejor el campo.''
\end{loreblock}

\vspace{0.4cm}

\section{Antes de la Batalla: Preparación}

El inicio de cada partida sigue una secuencia clara. Respetarla es parte de la táctica.

\begin{reglaespecial}{Secuencia de Preparación}
\begin{enumerate}
    \item \textbf{Definir el Escenario:} Se baraja el mazo de Misiones Principales y se roba la carta superior. Esta carta es pública y determina el mapa, las zonas de despliegue de cada facción y la condición de victoria principal.
    \item \textbf{Misiones Secundarias:} Cada jugador roba en secreto \textbf{2 cartas} del mazo de Misiones Secundarias. Estas son información privada.
    \item \textbf{Tirada de Iniciativa:} Cada jugador tira \textbf{1d12}. El ganador elige quién despliega primero.
    \item \textbf{Despliegue Alternado:} El primer jugador coloca una unidad en su zona de inicio (según la carta de Misión Principal). Luego el oponente coloca una unidad. Se alternan hasta que ambos jugadores hayan desplegado todo su ejército.
    \item \textbf{Inicio de la Primera Ronda.}
\end{enumerate}
\end{reglaespecial}

\section{Misión Principal}

La carta de Misión Principal define el objetivo público de la partida. Estas cartas tienen formato grande (similar a las cartas de Magic: The Gathering) e incluyen:

\begin{itemize}
    \item Un \textbf{mapa de zonas de inicio} que indica dónde puede desplegar cada facción.
    \item La \textbf{condición de victoria} principal (ver ejemplos abajo).
    \item Reglas especiales del escenario, si las hubiera.
\end{itemize}

\begin{imagebox}{Ejemplo de carta de Misión Principal — formato grande con mapa de zonas de despliegue impreso}
    \vspace{3cm}
\end{imagebox}

\subsection*{Ejemplos de Misiones Principales}

\begin{tabularx}{\linewidth}{|l|X|}
\hline
\rowcolor{HammerRed!15} \textbf{Nombre} & \textbf{Condición de Victoria} \\ \hline
\textbf{Aniquilación} & Elimina a todas las unidades enemigas. \\ \hline
\textbf{Robo de Gema} & Toma la Gema enemiga y deposítala en tu base. \\ \hline
\textbf{Resistencia} & El Equipo A defiende el centro por X rondas. El Equipo B debe romper la línea. \\ \hline
\textbf{Puntos de Control} & Ocupa y mantén la mayoría de los puntos de control al final de la ronda X. \\ \hline
\textbf{Caza de Reliquias} & Ocupa puntos de control por 1 ronda entera para robar una carta de Reliquia aleatoria del mazo. Gana quien acumule más puntos al final. \\ \hline
\end{tabularx}

\vspace{0.3cm}

\section{Misiones Secundarias (Secretas)}

Al inicio de la partida, cada jugador roba \textbf{2 cartas} del mazo de Misiones Secundarias y las mantiene en secreto. Estas misiones otorgan \textbf{Puntos de Victoria adicionales} al cumplirse.

\begin{reglaespecial}{Reglas de las Misiones Secundarias}
\begin{itemize}
    \item Las cartas son \textbf{información privada}. No se revelan al oponente hasta cumplirse.
    \item Al \textbf{comienzo de cada ronda}, cada jugador debe tener exactamente \textbf{2 cartas} en mano. Si una fue completada y descartada, se roba una nueva al iniciar la ronda siguiente.
    \item \textbf{Una vez por ronda}, un jugador puede descartar 1 carta de misión secundaria y robar 1 nueva en su lugar. Esta acción se realiza en la fase de inicio de ronda.
    \item Cuando se cumple una misión, se \textbf{revela la carta}, se anotan los Puntos de Victoria correspondientes, y la carta pasa al descarte.
\end{itemize}
\end{reglaespecial}

\subsection*{Ejemplos de Misiones Secundarias}

\begin{itemize}
    \item \textit{Territorio Seco:} Crea 3 Zonas Secas antes del final de la ronda 4.
    \item \textit{Línea Defensiva:} Mantén al menos 2 unidades en el cuadrante trasero propio durante 2 rondas consecutivas.
    \item \textit{Avance Imparable:} Mueve una unidad hasta el cuadrante trasero enemigo.
    \item \textit{Buscador de Reliquias:} Ocupa un punto de control durante 1 ronda entera para robar una carta de Reliquia.
    \item \textit{Retroceso Forzado:} Elimina todas las unidades enemigas de la zona central del mapa.
    \item \textit{Cacería Específica:} Elimina al Líder o Héroe enemigo.
\end{itemize}

\begin{imagebox}{Un jugador revelando una misión secundaria completada — cara oculta hasta el momento de la revelación}
    \vspace{2cm}
\end{imagebox}

\section{Despliegue de Unidades}

El despliegue es alternado y comienza una vez definida la Misión Principal.

\begin{reglaespecial}{El Despliegue Alternado}
\begin{enumerate}
    \item El jugador con iniciativa coloca \textbf{una unidad} dentro de su zona de inicio marcada en la carta de misión.
    \item El jugador contrario coloca \textbf{una unidad} en su zona.
    \item Se repite hasta que ambos ejércitos estén desplegados por completo.
    \item Las unidades \textbf{no pueden activarse} durante la fase de despliegue. El posicionamiento es la única acción disponible.
\end{enumerate}
\end{reglaespecial}

\subsection*{Unidades Caídas en el Tablero}

Cuando una unidad es eliminada en combate, su miniatura \textbf{se recuesta sobre el tablero} en el lugar donde cayó. Los cuerpos:

\begin{itemize}
    \item \textbf{Bloquean el movimiento} como terreno difícil (cuestan 1 PA extra para atravesar).
    \item Pueden ser objetivo de ciertas habilidades especiales de facción (necromancia, absorción de energía, etc.).
    \item No son removidos del tablero a menos que una regla especial lo indique.
\end{itemize}

\begin{tcolorbox}[colback=GoldRule!10, colframe=GoldRule, title=\textbf{REGLA: Terreno y Obstáculos}]
\small Los cuerpos, estructuras destruidas y Zonas Secas modifican el campo de batalla durante la partida. Consulta las reglas de Terreno de tu facción en el Codex correspondiente.
\end{tcolorbox}
